\documentclass[12pt, a4paper]{report}

\usepackage[utf8]{inputenc}
\usepackage[T1]{fontenc}
\usepackage[magyar]{babel}
\usepackage[hidelinks]{hyperref}
\usepackage{graphicx}
\usepackage{geometry}
\usepackage{listings}
\usepackage{xcolor}
\usepackage{titlesec}
\usepackage{setspace}
\usepackage{tocbibind}

% Oldalbeállítások (Szakdolgozati sztenderd)
\geometry{
    top=2.5cm,
    bottom=2.5cm,
    left=3cm,
    right=2.5cm
}
\onehalfspacing

% Kódblokkok formázása
\definecolor{codegray}{rgb}{0.5,0.5,0.5}
\definecolor{backcolour}{rgb}{0.95,0.95,0.92}
\definecolor{keywordcolor}{rgb}{0.13, 0.29, 0.53}
\definecolor{stringcolor}{rgb}{0.31, 0.60, 0.02}

\lstdefinestyle{mystyle}{
    backgroundcolor=\color{backcolour},   
    commentstyle=\color{codegray}\itshape,
    keywordstyle=\color{keywordcolor}\bfseries,
    numberstyle=\tiny\color{codegray},
    stringstyle=\color{stringcolor},
    basicstyle=\ttfamily\footnotesize,
    breakatwhitespace=false,         
    breaklines=true,                 
    captionpos=b,                    
    keepspaces=true,                 
    numbers=left,                    
    numbersep=5pt,                  
    showspaces=false,                
    showstringspaces=false,
    showtabs=false,                  
    tabsize=2
}
\lstset{style=mystyle, language=HTML} % Alapértelmezett, de TSX-hez is jó lesz

\begin{document}

% --- CÍMLAP ---
\begin{titlepage}
    \centering
    \vspace*{3cm}
    
    {\scshape\Large Szakmai Gyakorlat és FFT Projekt\par}
    \vspace{1.5cm}
    
    {\huge\bfseries Színjátszó Tábor Weboldala\par}
    {\Large\bfseries (Csalamádé Next.js Projekt) Átfogó Dokumentáció\par}
    \vspace{2cm}
    
    {\Large Készítette:\\ \textbf{Nagy Csaba}}\par
    \vspace{0.5cm}
    {\large Számítástechnika IV.\par}
    
    \vfill
    
    % Bottom of the page
    {\large 2026\par}
\end{titlepage}

% --- TARTALOMJEGYZÉK ---
\tableofcontents
\newpage

% --- 1. BEVEZETÉS ---
\chapter{Bevezetés}

A modern webfejlesztés rohamos fejlődésével a statikus, tisztán információközlésre fókuszáló weboldalak ideje lejárt, különösen akkor, ha a célközönség a vizuális ingerekhez szokott Z és Alfa generáció. A jelen dokumentáció a "Csalamádé" néven ismert diákszínjátszó tábor interaktív weboldalának felépítését tárgyalja, a kezdeti tervezési fázisoktól (FFT) a komplex szoftvermérnöki architektúra felépítéséig.

A dokumentáció részletes technikai elemzést nyújt a projekt forráskódjáról, kitérve arra, hogy miért éppen a **Next.js**, a **TanStack Query** és a **Tailwind CSS** technológiákat választottuk, és hogyan valósítottuk meg a komplex, mobilon is kiválóan teljesítő komponenseket.

% --- 2. FELHASZNÁLÓI FELÜLETEK TERVEZÉSE (FFT) ---
\chapter{Felhasználói Felületek Tervezése (FFT)}

A fejlesztés nulladik lépése a potenciális látogatók azonosítása és a felhasználói élmény (UX) megtervezése volt.

\section{Célközönség és Personák}
A weboldal kettős célközönséggel rendelkezik, amelyeket a következő personák reprezentálnak:

\begin{itemize}
    \item \textbf{Persona 1: A táborozó diák (16-18 éves):} Olyan vizuális ingerekre vágyik, amelyek megmutatják, milyen lesz a tábor hangulata, illetve kik lesznek ott. Egy unalmas, elavult kinézetű oldal elriasztja. Számára az élmény és a videós tartalom a legfontosabb.
    \item \textbf{Persona 2: Nagy Gábor (44 éves, szülő):} Olyan biztonságos, fejlesztő környezetet keres a lányának, ahol értékteremtő munka folyik szakemberek vezetésével. Fájdalompontja, hogy sok tábornak nincs hivatalos weboldala, átláthatatlanok a szervezők, és zavarja a bonyolult, papíralapú jelentkezési formaság.
\end{itemize}

\section{Piackutatás és Felhasználói Vélemények}
A tervezést egy átfogó piackutatás előzte meg, mely során hasonló táborok online jelenlétét vizsgáltuk:
\begin{itemize}
    \item \textbf{Hagyományos tábori oldalak (pl. WordPress):} Előnyük, hogy az információ pontos, de hátrányuk a statikus, szövegközpontú elrendezés és a lassú, elavult mobilnézet, amely nem rezonál a fiatalabb generációval.
    \item \textbf{Kizárólag Facebook Page használata:} Jó az elérés, de az információk elvesznek a hírfolyamban. Nincs dedikált jelentkezési felület és nehezen kereshetők vissza az archív anyagok.
\end{itemize}

Kérdőíves felméréseink és mélyinterjúk során kiderült, hogy a diákok 80\%-a egy-egy hangulatvideó alapján dönti el, érdekli-e a rendezvény. Ezek alapján született meg a mi \textbf{új, kiterjesztett ötletünk}: Információközlés helyett \textit{immerzív élményt} nyújtunk dinamikus hátterekkel, interaktív videós kártyákkal és mozgásközpontú animációkkal.

\section{Vizuális Koncepció (Design Tokenek)}
A Figma tervek alapján a végső koncepció a következő vizuális irányelveket követi:
\begin{itemize}
    \item \textbf{Színek és Téma:} A sötét mód dominál. Ez szükséges a "Pop-out" hatás eléréséhez, így a színes videó-bélyegképek kitűnnek a háttérből.
    \item \textbf{Tipográfia:} A \textit{Museo} betűcsalád (Museo300 és Museo700) használata, amely kerekded, mégis komoly karaktere miatt tökéletesen passzol a játékos, színházi környezethez.
\end{itemize}

% --- 3. FRONTEND ARCHITEKTÚRA ÉS NEXT.JS ---
\chapter{Frontend Architektúra és Keretrendszerek}

A vizuális tervek elkészülte után a legfontosabb döntés a megfelelő technológiai stack kiválasztása volt. A választás a **Next.js** App Router alapú architektúrájára esett. 

\section{Miért a Next.js?}
A klasszikus React SPA (Single Page Application) legnagyobb hátránya a gyenge keresőoptimalizálás (SEO) és a lassú kezdeti oldalbetöltés, mivel a böngészőnek le kell töltenie a teljes JavaScript bundle-t a renderelés megkezdése előtt. A Next.js **SSR (Server-Side Rendering)** képessége megoldja ezt a problémát: a szerver már a kész HTML-t küldi a kliensnek, így az oldal azonnal megjelenik, és a keresőrobotok (Google) is tökéletesen tudják indexelni az eseményt.

\section{A Root Layout Felépítése (layout.tsx)}
Minden Next.js (App Router) projekt lelke a \texttt{layout.tsx}. Itt konfiguráljuk a globális szolgáltatókat (Providereket) és töltjük be a betűtípusokat a legoptimálisabb módon a \texttt{next/font} API használatával.

\begin{lstlisting}[language=HTML, caption=A projekt alapját képező app/layout.tsx részlete]
import { LanguageProvider } from "@/context/LanguageContext";
import QueryProvider from "@/components/QueryProvider";
import localFont from "next/font/local";

// Helyi betutipusok betoltese es CSS valtozokkent valo regisztralasa
const museo300 = localFont({
  src: "../public/fonts/Museo300.otf",
  variable: "--font-museo-300",
  display: "swap", // Optimalizalt fontmegjelenites
});

export default function RootLayout({ children }: { children: React.ReactNode }) {
  return (
    <html lang="hu" suppressHydrationWarning>
      <body className={`${museo300.variable} antialiased`}>
        {/* Adatkezelesi reteg */}
        <QueryProvider>
          {/* Nyelvi kontextus */}
          <LanguageProvider>
            <main>{children}</main>
          </LanguageProvider>
        </QueryProvider>
      </body>
    </html>
  );
}
\end{lstlisting}

**Mit csinál a fenti kód?**
1. **\texttt{localFont}**: Beolvassa a lokális \textit{Museo} fontfájlokat, optimalizálja azokat, és egyedi CSS változóként (\texttt{--font-museo-300}) adja át a DOM-nak. Ez megakadályozza a CLS (Cumulative Layout Shift) jelenséget, ami drasztikusan javítja a Google Lighthouse pontszámot.
2. **\texttt{QueryProvider} \& \texttt{LanguageProvider}**: Ezek az úgynevezett Context Providerek. A Next.js 13+ App Routerében alapértelmezetten minden komponens Server Component (szerveren fut le). Mivel a nyelvi váltás és az adatlekérdezés gyorsítótára (cache) kliens-oldali interaktivitást igényel, ezekbe a providerekbe csomagoljuk az alkalmazást.

% --- 4. ÁLLAPOTKEZELÉS ÉS ADATÁRAMLÁS ---
\chapter{Állapotkezelés: TanStack Query}

Egy modern webalkalmazás szinte sosincs elszigetelve; adatokat kell lekérnie API-kon keresztül (pl. videók adatai, résztvevők listája). Ezt a klasszikus Reactban \texttt{useEffect} és \texttt{useState} segítségével szokás megoldani, amely azonban rendkívül sebezhető: nincs beépített cache, nincs automatikus újrapróbálkozás hiba esetén, és a komponens újrarajzolásakor felesleges hálózati kérések indulnak.

A projektben erre a \textbf{TanStack Query} (korábban React Query) technológiát használtuk.

\section{A QueryProvider Működése}

\begin{lstlisting}[language=HTML, caption=A QueryProvider.tsx implementációja]
"use client";

import { QueryClient, QueryClientProvider } from '@tanstack/react-query';
import { useState } from 'react';

export default function QueryProvider({ children }: { children: React.ReactNode }) {
  const [queryClient] = useState(
    () =>
      new QueryClient({
        defaultOptions: {
          queries: {
            staleTime: 60 * 1000, // 1 percig frissnek tekinti az adatot
            retry: 1, // Hiba eseten egyszer probalkozik ujra
          },
        },
      })
  );

  return (
    <QueryClientProvider client={queryClient}>
      {children}
    </QueryClientProvider>
  );
}
\end{lstlisting}

Ebben a kódrészletben a \texttt{staleTime: 60 * 1000} beállítás garantálja, hogy ha a felhasználó navigál az oldalon és visszatér egy olyan szekcióhoz (pl. a Táborok listájához), az alkalmazás nem küld új API kérést, hanem azonnal, 0 milliszekundum alatt kiszolgálja a gyorsítótárból a memóriában lévő adatot. Ezzel hatalmas szerverterhelést spórolunk meg, és a felhasználói élmény hihetetlenül fluid marad.

% --- 5. DINAMIKUS KOMPONENSEK: CAMPS SECTION ---
\chapter{Dinamikus Komponensek és Reszponzivitás}

A projekt legkomplexebb és mérnökileg legnagyobb kihívást jelentő része a \texttt{CampsSection.tsx} volt. Ennek a szekciónak a feladata a korábbi évek táborainak (videók) megjelenítése egy interaktív, vízszintesen görgethető (Carousel) formátumban.

\section{A probléma: Mobilos navigáció és a "Two-Tap" logika}
Amikor a kártyák vízszintesen vannak felsorakoztatva, a felhasználó ujjal (swipe) próbálja görgetni azokat okostelefonon. A fejlesztés korai szakaszában a swipe mozdulat véletlenül rákattintott egy kártyára, ami váratlanul átirányította a felhasználót a YouTube-ra. Ez rendkívül frusztráló UX-et eredményezett.

\textbf{A mérnöki megoldás:} A "Two-Tap" (dupla koppintásos) logika implementálása aszinkron fókuszkövetéssel.

\begin{lstlisting}[language=HTML, caption=Részlet a CampsSection.tsx kódból (Intersection Observer)]
useEffect(() => {
    if (!config.isPortrait) return; // Csak mobilon fut le
    
    const timeoutId = setTimeout(() => {
        const observer = new IntersectionObserver(
            (entries) => {
                entries.forEach((entry) => {
                    if (entry.isIntersecting) {
                        const cardId = parseInt(entry.target.getAttribute('data-card-id') || '1');
                        setCenteredCardId((prev) => {
                            if (prev !== cardId) {
                                setTappedCardId(null); // Uj focusnal elveszitjuk a kijelolest
                                return cardId;
                            }
                            return prev;
                        });
                    }
                });
            },
            {
                root: scrollContainerRef.current,
                rootMargin: '0px -45% 0px -45%', // Pontosan a kepernyo kozepet vizsgalja
                threshold: 0.1,
            }
        );

        cardRefs.current.forEach((element) => {
            if (element) observer.observe(element);
        });

        return () => observer.disconnect(); // Takaritas
    }, 100);

    return () => clearTimeout(timeoutId);
}, [config.isPortrait]);
\end{lstlisting}

**Mit csinál a kód?**
Ahelyett, hogy drága (számításigényes) görgetési eseményeket (\texttt{onScroll}) figyelnénk, a böngésző natív, rendkívül hardvergyorsított \textbf{Intersection Observer API}-ját használjuk. 
A \texttt{rootMargin: '0px -45\% 0px -45\%'} beállítás gyakorlatilag egy apró, láthatatlan vonalat hoz létre a képernyő közepén. Amikor a felhasználó görget egy kártyát, és az metszi ezt a középső vonalat, a kártya bekerül a \texttt{centeredCardId} állapotba (state). Ekkor színes lesz és megnő. A \textit{második} tudatos érintésre (\texttt{isTapped}) pedig elindul a külső videós link. Ezzel a frusztráló kattintási hibákat 100\%-ban elimináltuk.

\section{A Reszponzív Állapotgép (State Machine)}
A komponens folyamatosan figyeli a képernyőméretet a \texttt{performResize} függvény segítségével (Debouncing technikával, hogy ne terhelje a processzort átméretezéskor). Mobilon dinamikusan változtatja a kártyák méretét (\texttt{vw} és \texttt{vh} egységek használatával), míg asztali gépen egy mágneses (hover) interakciót tesz lehetővé.

% --- 6. UI/UX ÉS ANIMÁCIÓK ---
\chapter{UI/UX Implementáció és Animációk (GSAP)}

Az oldal dizájnja szándékosan szakít a statikus sablonokkal. A vizuális visszajelzések és a fluiditás a márka (Csalamádé) alapértéke.

\section{Tailwind CSS a gyorsaságért}
A komponensek megírásánál a Tailwind CSS-t használtuk. Egy klasszikus BEM architektúrájú CSS helyett a utility-first osztályok (pl. \texttt{flex flex-col items-center justify-center pt-0 -mt-30 px-0 pb-10 relative z-10}) lehetővé tették, hogy a stílus a komponens logikájával egy fájlban (\textit{Colocation}) maradjon. Ez javítja a fejlesztési élményt és csökkenti a halott CSS kód mennyiségét, hiszen a fordító (compiler) csak a ténylegesen használt osztályokat építi be a végső CSS fájlba.

\section{GSAP és Mikro-interakciók}
A React beépített animációs könyvtárai (pl. Framer Motion) helyett a projekt a \textbf{GSAP (GreenSock Animation Platform)} technológiára épít a nagyobb teljesítmény és kompatibilitás érdekében. A GSAP képes közvetlenül a böngésző GPU-jára küldeni a transzformációs utasításokat. A \texttt{/jelentkezz} végponton futó SVG maszkolás, ahol a tábor logója "kitágul", és a felhasználó fizikailag beesik az űrlapba, egy klasszikus GSAP \textit{Timeline} segítségével készült, megszakítva a hagyományos "kattintok-és-betölt-egy-új-oldal" monotóniáját.

% --- 7. ÖSSZEGZÉS ---
\chapter{Összegzés és Szoftvermérnöki Tanulságok}

A Csalamádé Színjátszó Tábor weboldalának fejlesztése kiváló lehetőséget biztosított arra, hogy a szoftverfejlesztési és objektumorientált elveket egy modern, piacképes termékben integráljuk.

\section{Tanulságok és Kihívások}
A legnagyobb kihívást az animációk, a teljesítmény és az akadálymentesség egyensúlyának megtartása jelentette. Míg az asztali böngészők hardveres gyorsítással könnyen kezelik a GSAP animációkat és a nagyfelbontású \texttt{next/image} képeket, addig az olcsóbb, régebbi okostelefonokon agresszív memória-optimalizációra (WebP formátum, csökkentett polyfill) volt szükség a zökkenőmentes görgetés érdekében. Az iOS Safari különleges DOM-kezelési sajátosságai (például a z-index hibák a 3D transzformációknál) számos egyedi fallback írását tették szükségessé.

\section{További fejlesztési lehetőségek (Future Work)}
Bár a jelenlegi verzió rendkívül robusztus és esztétikus, a rendszer jövőbeni skálázhatósága az alábbi modulokkal bővülhet:
\begin{enumerate}
    \item \textbf{Integrált Jelentkezési Backend:} A külső Google Forms leváltása, egy saját Next.js API Route alapú űrlapkezelőre, amely közvetlenül egy relációs adatbázisba (pl. PostgreSQL) dolgozik.
    \item \textbf{Headless CMS Integráció:} Egy tartalomkezelő rendszer (pl. Strapi vagy Sanity) bekötése, hogy a tábor szervezői programozói tudás nélkül tudják cserélni a videókat és az idézeteket.
\end{enumerate}

Összességében a projekt maradéktalanul teljesítette a célkitűzéseket: a hagyományos, statikus információmegosztást egy dinamikus, vizuálisan magával ragadó webes élménnyé transzformálta.

\end{document}
